% !TeX root = part2.tex

\documentclass[../final.tex]{subfiles}

\begin{document}

% ============================================================
% Семинар 2
% ============================================================

\practice{2}{12.09.2026}{Колебания решётки}


% ============================================================
% Задача 1
% ============================================================

\task[1]{Линейная цепочка с одним атомом в элементарной ячейке}

\condition
Определить собственные моды, то есть фононный спектр,
бесконечной одномерной цепочки, состоящей из одинаковых
атомов массы $m$, соединённых одинаковыми пружинками
жёсткости $\alpha$.


% ------------------------------------------------------------
% РИСУНОК 1
% ------------------------------------------------------------

\begin{rbox}[left=18pt,right=18pt,top=14pt,bottom=10pt]{[]}
    \centering

    \begin{tikzpicture}[
        >=Latex,
        every node/.style={text=rtext}
    ]

        % Пружины
        \foreach \x in {-4,-2,0,2}{
            \draw[
                draw=rtext,
                line width=0.8pt
            ]
            plot[
                domain=0:1,
                samples=70,
                variable=\t
            ]
            (
                {\x+0.28+1.44*\t},
                {0.11*sin(1440*\t)}
            );
        }

        % Продолжение цепочки
        \node at (-5.0,0) {$\ldots$};
        \node at (5.0,0) {$\ldots$};

        % Атомы
        \foreach \x in {-4,-2,2,4}{
            \filldraw[
                fill=rbgsecondary,
                draw=rtext,
                line width=0.9pt
            ]
            (\x,0) circle (0.28);

            \node[above=5pt] at (\x,0) {$m$};
        }

        % n-й атом
        \filldraw[
            fill=rc3,
            draw=rtext,
            line width=0.9pt
        ]
        (0,0) circle (0.28);

        \node[above=5pt] at (0,0) {$m$};
        \node[below=5pt] at (0,0) {$n$};

        % Жёсткости
        \node[above] at (-3,0.18) {$\alpha$};
        \node[above] at (-1,0.18) {$\alpha$};
        \node[above] at (1,0.18) {$\alpha$};
        \node[above] at (3,0.18) {$\alpha$};

        % Период решётки
        \draw[
            <->,
            draw=rtext,
            line width=0.8pt
        ]
        (-2,-0.85)
        --
        (0,-0.85)
        node[midway,below] {$a$};

        % Смещение
        \draw[
            ->,
            draw=rc3,
            line width=1.0pt
        ]
        (0,0.55)
        --
        (0.9,0.55)
        node[midway,above] {$u_n$};

    \end{tikzpicture}

    \captionsetup{
        font=footnotesize,
        labelfont=bf,
        skip=5pt
    }

    \captionof{figure}{
        Одномерная цепочка одинаковых атомов массы $m$,
        соединённых пружинами жёсткости $\alpha$.
        Расстояние между соседними положениями равновесия равно $a$,
        а $u_n$ обозначает смещение $n$-го атома.
    }

    \label{fig:monoatomic_chain}
\end{rbox}


Рассмотрим одномерную цепочку одинаковых атомов.
Положение равновесия $n$-го атома можно записать как
%
\begin{equation*}
    x_n=na,
\end{equation*}
%
где $a$ --- период решётки.

Пусть
%
\begin{equation*}
    u_n(t)
\end{equation*}
%
--- смещение $n$-го атома от положения равновесия.


% ------------------------------------------------------------
% Лагранжиан
% ------------------------------------------------------------

\subsection*{Лагранжиан цепочки}

Кинетическая энергия цепочки равна
%
\begin{equation*}
    T
    =
    \sum_n
    \frac{m}{2}
    \dot u_n^{\,2}.
\end{equation*}

При малых колебаниях каждую связь между соседними атомами
считаем гармонической.

Удлинение пружины между атомами $n$ и $n+1$ равно
%
\begin{equation*}
    u_{n+1}-u_n.
\end{equation*}

Поэтому потенциальная энергия этой пружины равна
%
\begin{equation*}
    \frac{\alpha}{2}
    (u_{n+1}-u_n)^2.
\end{equation*}

Суммарная потенциальная энергия цепочки имеет вид
%
\begin{equation*}
    U
    =
    \sum_n
    \frac{\alpha}{2}
    (u_n-u_{n+1})^2.
\end{equation*}

Следовательно,
%
\begin{equation*}
    L=T-U,
\end{equation*}
%
то есть
%
\begin{equation*}
    \rboxed{
        L
        =
        \sum_n
        \left[
            \frac{m}{2}\dot u_n^{\,2}
            -
            \frac{\alpha}{2}
            (u_n-u_{n+1})^2
        \right]
    }.
\end{equation*}


% ------------------------------------------------------------
% Уравнение движения
% ------------------------------------------------------------

\subsection*{Уравнение движения}

Для каждой координаты $u_n$ записываем
уравнение Эйлера--Лагранжа:
%
\begin{equation*}
    \frac{d}{dt}
    \frac{\partial L}{\partial\dot u_n}
    -
    \frac{\partial L}{\partial u_n}
    =
    0.
\end{equation*}

Сначала вычислим производную по скорости:
%
\begin{equation*}
    \frac{\partial L}{\partial\dot u_n}
    =
    m\dot u_n,
\end{equation*}
%
поэтому
%
\begin{equation*}
    \frac{d}{dt}
    \frac{\partial L}{\partial\dot u_n}
    =
    m\ddot u_n.
\end{equation*}

Теперь найдём
%
\begin{equation*}
    \frac{\partial L}{\partial u_n}.
\end{equation*}

Координата $u_n$ входит сразу в два соседних упругих члена:
%
\begin{equation*}
    -\frac{\alpha}{2}(u_n-u_{n+1})^2,
    \qquad
    -\frac{\alpha}{2}(u_{n-1}-u_n)^2.
\end{equation*}

Поэтому
%
\begin{align*}
    \frac{\partial L}{\partial u_n}
    &=
    -\alpha(u_n-u_{n+1})
    +
    \alpha(u_{n-1}-u_n)
    \\
    &=
    \alpha
    \left(
        u_{n+1}
        +
        u_{n-1}
        -
        2u_n
    \right).
\end{align*}

Таким образом,
%
\begin{equation*}
    \rboxed{
        m\ddot u_n
        =
        \alpha
        \left(
            u_{n+1}
            +
            u_{n-1}
            -
            2u_n
        \right)
    }.
\end{equation*}


% ------------------------------------------------------------
% Физический смысл силы
% ------------------------------------------------------------

\subsection*{Сила, действующая на атом}

Правая часть уравнения движения представляет собой
результирующую силу двух соседних пружин.

Сила со стороны правой пружины:
%
\begin{equation*}
    F_{n,n+1}
    =
    \alpha(u_{n+1}-u_n).
\end{equation*}

Сила со стороны левой пружины:
%
\begin{equation*}
    F_{n,n-1}
    =
    \alpha(u_{n-1}-u_n).
\end{equation*}

Поэтому
%
\begin{align*}
    F_n
    &=
    F_{n,n+1}
    +
    F_{n,n-1}
    \\
    &=
    \alpha(u_{n+1}-u_n)
    +
    \alpha(u_{n-1}-u_n)
    \\
    &=
    \alpha
    \left(
        u_{n+1}
        +
        u_{n-1}
        -
        2u_n
    \right).
\end{align*}

Следовательно,
%
\begin{equation*}
    m\ddot u_n
    =
    \alpha
    \left(
        u_{n+1}
        +
        u_{n-1}
        -
        2u_n
    \right).
\end{equation*}


% ------------------------------------------------------------
% Нормальная мода
% ------------------------------------------------------------

\subsection*{Нормальная мода}

Ищем нормальное колебание в виде плоской волны:
%
\begin{equation*}
    \rboxed{
        u_n
        =
        u
        e^{-i\omega t+ikna}
    }.
\end{equation*}

Здесь $k$ --- волновой вектор колебания,
$\omega$ --- частота нормальной моды,
а $a$ --- период цепочки.

Для соседних атомов имеем
%
\begin{equation*}
    u_{n+1}
    =
    u_n e^{ika},
    \qquad
    u_{n-1}
    =
    u_n e^{-ika}.
\end{equation*}

Кроме того,
%
\begin{equation*}
    \ddot u_n
    =
    -\omega^2u_n.
\end{equation*}

Подставляем всё в уравнение движения:
%
\begin{equation*}
    -m\omega^2u_n
    =
    \alpha u_n
    \left(
        e^{ika}
        +
        e^{-ika}
        -
        2
    \right).
\end{equation*}

Сокращаем на ненулевую амплитуду $u_n$:
%
\begin{equation*}
    -m\omega^2
    =
    \alpha
    \left(
        e^{ika}
        +
        e^{-ika}
        -
        2
    \right).
\end{equation*}

Используем
%
\begin{equation*}
    e^{ika}+e^{-ika}
    =
    2\cos ka.
\end{equation*}

Тогда
%
\begin{equation*}
    -m\omega^2
    =
    2\alpha
    \left(
        \cos ka-1
    \right),
\end{equation*}
%
или
%
\begin{equation*}
    m\omega^2
    =
    2\alpha
    \left(
        1-\cos ka
    \right).
\end{equation*}

Используем тождество
%
\begin{equation*}
    1-\cos x
    =
    2\sin^2\frac{x}{2}.
\end{equation*}

Получаем
%
\begin{equation*}
    \rboxed{
        \omega^2(k)
        =
        \frac{4\alpha}{m}
        \sin^2\frac{ka}{2}
    }.
\end{equation*}

Следовательно,
%
\begin{equation*}
    \rboxed{
        \omega(k)
        =
        2\sqrt{\frac{\alpha}{m}}
        \left|
            \sin\frac{ka}{2}
        \right|
    }.
\end{equation*}


% ------------------------------------------------------------
% Периодичность спектра
% ------------------------------------------------------------

\subsection*{Периодичность спектра}

Поскольку
%
\begin{equation*}
    \sin^2
    \left(
        \frac{ka}{2}+\pi
    \right)
    =
    \sin^2\frac{ka}{2},
\end{equation*}
%
спектр периодичен по $k$:
%
\begin{equation*}
    \rboxed{
        \omega
        \left(
            k+\frac{2\pi}{a}
        \right)
        =
        \omega(k)
    }.
\end{equation*}

Поэтому достаточно рассматривать первую зону Бриллюэна:
%
\begin{equation*}
    \rboxed{
        -\frac{\pi}{a}
        \leq
        k
        \leq
        \frac{\pi}{a}
    }.
\end{equation*}

В центре зоны,
%
\begin{equation*}
    k=0,
\end{equation*}
%
имеем
%
\begin{equation*}
    \omega(0)=0.
\end{equation*}

На границе зоны,
%
\begin{equation*}
    k=\pm\frac{\pi}{a},
\end{equation*}
%
частота достигает максимального значения:
%
\begin{equation*}
    \rboxed{
        \omega_{\max}
        =
        2\sqrt{\frac{\alpha}{m}}
    }.
\end{equation*}

Или
%
\begin{equation*}
    \rboxed{
        \omega_{\max}^{\,2}
        =
        \frac{4\alpha}{m}
    }.
\end{equation*}


% ------------------------------------------------------------
% РИСУНОК 2
% ------------------------------------------------------------

\begin{rbox}[left=18pt,right=18pt,top=14pt,bottom=10pt]{[]}
    \centering

    \begin{tikzpicture}[
        x=0.72cm,
        y=1.45cm,
        >=Latex,
        every node/.style={text=rtext}
    ]

        % Оси
        \draw[
            ->,
            draw=rtext,
            line width=0.9pt
        ]
        (-6.8,0)
        --
        (6.9,0)
        node[right] {$k$};

        \draw[
            ->,
            draw=rtext,
            line width=0.9pt
        ]
        (0,0)
        --
        (0,2.25)
        node[above] {$\omega$};

        % Первая зона Бриллюэна
        \draw[
            draw=rtext,
            dashed,
            opacity=0.65
        ]
        (-3.1416,0)
        --
        (-3.1416,1.95);

        \draw[
            draw=rtext,
            dashed,
            opacity=0.65
        ]
        (3.1416,0)
        --
        (3.1416,1.95);

        % Дисперсионная кривая
        \draw[
            draw=rc3,
            line width=1.2pt,
            domain=-6.2832:6.2832,
            samples=250,
            variable=\x
        ]
        plot
        (
            {\x},
            {1.75*abs(sin(deg(\x/2)))}
        );

        % Деления
        \foreach \x/\lab in {
            -6.2832/{-\frac{2\pi}{a}},
            -3.1416/{-\frac{\pi}{a}},
             3.1416/{\frac{\pi}{a}},
             6.2832/{\frac{2\pi}{a}}
        }{
            \draw[
                draw=rtext,
                line width=0.7pt
            ]
            (\x,0.06)
            --
            (\x,-0.06);

            \node[below=3pt] at (\x,0) {$\lab$};
        }

        % Максимальная частота
        \draw[
            draw=rtext,
            dashed,
            opacity=0.55
        ]
        (-6.2,1.75)
        --
        (6.2,1.75);

        \node[
            left
        ]
        at (0,1.75)
        {$2\sqrt{\alpha/m}$};

    \end{tikzpicture}

    \captionsetup{
        font=footnotesize,
        labelfont=bf,
        skip=5pt
    }

    \captionof{figure}{
        Дисперсионная зависимость
        $\omega(k)=2\sqrt{\alpha/m}\,
        |\sin(ka/2)|$
        для одноатомной линейной цепочки.
        Пунктиром отмечены границы первой зоны Бриллюэна.
    }

    \label{fig:monoatomic_dispersion}
\end{rbox}


Таким образом, спектр распространяющихся колебаний ограничен:
%
\begin{equation*}
    \rboxed{
        0
        \leq
        \omega
        \leq
        2\sqrt{\frac{\alpha}{m}}
    }.
\end{equation*}

Частоты
%
\begin{equation*}
    \omega>\omega_{\max}
\end{equation*}
%
не соответствуют распространяющимся нормальным модам
идеальной цепочки.


% ============================================================
% Задача 2
% ============================================================

\task[2]{Примесный атом в линейной цепочке: другая масса}

\condition
Масса $M$ одного из атомов цепочки отличается
от массы остальных атомов:
%
\begin{equation*}
    M\neq m.
\end{equation*}
%
Определить, как наличие примеси влияет на фононный спектр,
и найти условие существования локализованного состояния.


% ------------------------------------------------------------
% РИСУНОК 3
% ------------------------------------------------------------

\begin{rbox}[left=18pt,right=18pt,top=14pt,bottom=10pt]{[]}
    \centering

    \begin{tikzpicture}[
        >=Latex,
        every node/.style={text=rtext}
    ]

        % Пружины
        \foreach \x in {-4,-2,0,2}{
            \draw[
                draw=rtext,
                line width=0.8pt
            ]
            plot[
                domain=0:1,
                samples=70,
                variable=\t
            ]
            (
                {\x+0.28+1.44*\t},
                {0.11*sin(1440*\t)}
            );
        }

        % Края цепочки
        \node at (-5.0,0) {$\ldots$};
        \node at (5.0,0) {$\ldots$};

        % Обычные атомы
        \foreach \x in {-4,-2,2,4}{
            \filldraw[
                fill=rbgsecondary,
                draw=rtext,
                line width=0.9pt
            ]
            (\x,0) circle (0.28);

            \node[above=5pt] at (\x,0) {$m$};
        }

        % Примесь
        \filldraw[
            fill=rc3,
            draw=rtext,
            line width=1.0pt
        ]
        (0,0) circle (0.34);

        \node[above=6pt] at (0,0) {$M$};
        \node[below=6pt] at (0,0) {$n=0$};

        % Жёсткости
        \node[above] at (-3,0.18) {$\alpha$};
        \node[above] at (-1,0.18) {$\alpha$};
        \node[above] at (1,0.18) {$\alpha$};
        \node[above] at (3,0.18) {$\alpha$};

    \end{tikzpicture}

    \captionsetup{
        font=footnotesize,
        labelfont=bf,
        skip=5pt
    }

    \captionof{figure}{
        Одномерная цепочка с примесным атомом массы $M$.
        Все остальные атомы имеют массу $m$,
        а жёсткости всех пружин одинаковы и равны $\alpha$.
    }

    \label{fig:mass_impurity}
\end{rbox}


Поместим примесный атом в узел
%
\begin{equation*}
    n=0.
\end{equation*}

Для всех обычных атомов,
%
\begin{equation*}
    n\neq0,
\end{equation*}
%
уравнение движения имеет прежний вид:
%
\begin{equation*}
    \rboxed{
        m\ddot u_n
        =
        \alpha
        \left(
            u_{n+1}
            +
            u_{n-1}
            -
            2u_n
        \right)
    }.
\end{equation*}

Для примесного атома получаем
%
\begin{equation*}
    \rboxed{
        M\ddot u_0
        =
        \alpha
        \left(
            u_1
            +
            u_{-1}
            -
            2u_0
        \right)
    }.
\end{equation*}


% ------------------------------------------------------------
% Локализованная мода
% ------------------------------------------------------------

\subsection*{Локализованная мода}

Для локализованного состояния амплитуда колебаний должна
уменьшаться при удалении от примеси.

Ищем решение в виде
%
\begin{equation*}
    \rboxed{
        u_n
        =
        u_0
        (-1)^n
        e^{-\chi a|n|}
        e^{-i\omega t}
    },
    \qquad
    \chi>0.
\end{equation*}

Множитель
%
\begin{equation*}
    e^{-\chi a|n|}
\end{equation*}
%
отвечает за экспоненциальное затухание амплитуды
при удалении от примесного атома.

Множитель
%
\begin{equation*}
    (-1)^n
\end{equation*}
%
означает, что соседние атомы колеблются в противофазе.

Такое решение можно представить как продолжение обычной
плоской волны в область комплексных $k$:
%
\begin{equation*}
    k
    =
    \frac{\pi}{a}
    +
    i\chi.
\end{equation*}

Действительно,
%
\begin{align*}
    e^{ikna}
    &=
    e^{i(\pi/a+i\chi)na}
    \\
    &=
    e^{i\pi n}
    e^{-\chi na}
    \\
    &=
    (-1)^n
    e^{-\chi na}.
\end{align*}

Именно поэтому для локализованной моды появляется
дополнительный множитель $(-1)^n$.


% ------------------------------------------------------------
% Уравнение для обычных атомов
% ------------------------------------------------------------

\subsection*{Уравнение вдали от примеси}

Рассмотрим сначала область
%
\begin{equation*}
    n>0.
\end{equation*}

Тогда
%
\begin{equation*}
    u_{n+1}
    =
    -u_n e^{-\chi a},
    \qquad
    u_{n-1}
    =
    -u_n e^{\chi a}.
\end{equation*}

Кроме того,
%
\begin{equation*}
    \ddot u_n
    =
    -\omega^2u_n.
\end{equation*}

Подставляем в уравнение движения обычного атома:
%
\begin{equation*}
    -m\omega^2u_n
    =
    \alpha
    \left(
        -u_ne^{-\chi a}
        -
        u_ne^{\chi a}
        -
        2u_n
    \right).
\end{equation*}

После сокращения на $u_n$:
%
\begin{equation*}
    \rboxed{
        m\omega^2
        =
        \alpha
        \left(
            e^{\chi a}
            +
            e^{-\chi a}
            +
            2
        \right)
    }.
\end{equation*}

Аналогичное уравнение получается и для области $n<0$.


% ------------------------------------------------------------
% Уравнение для примесного атома
% ------------------------------------------------------------

\subsection*{Уравнение для примесного атома}

Для ближайших соседей примеси имеем
%
\begin{equation*}
    u_1
    =
    u_{-1}
    =
    -u_0e^{-\chi a}.
\end{equation*}

Подставляем это в уравнение для $n=0$:
%
\begin{equation*}
    -M\omega^2u_0
    =
    \alpha
    \left(
        -2u_0e^{-\chi a}
        -
        2u_0
    \right).
\end{equation*}

После сокращения:
%
\begin{equation*}
    \rboxed{
        M\omega^2
        =
        2\alpha
        \left(
            1+e^{-\chi a}
        \right)
    }.
\end{equation*}


% ------------------------------------------------------------
% Условие локализации
% ------------------------------------------------------------

\subsection*{Условие существования локализованного состояния}

Получены два выражения:
%
\begin{equation*}
    m\omega^2
    =
    \alpha
    \left(
        e^{\chi a}
        +
        e^{-\chi a}
        +
        2
    \right),
\end{equation*}
%
и
%
\begin{equation*}
    M\omega^2
    =
    2\alpha
    \left(
        1+e^{-\chi a}
    \right).
\end{equation*}

Делим второе равенство на первое:
%
\begin{equation*}
    \frac{M}{m}
    =
    \frac{
        2(1+e^{-\chi a})
    }{
        e^{\chi a}
        +
        e^{-\chi a}
        +
        2
    }.
\end{equation*}

Заметим, что
%
\begin{equation*}
    e^{\chi a}
    +
    e^{-\chi a}
    +
    2
    =
    \left(
        1+e^{\chi a}
    \right)
    \left(
        1+e^{-\chi a}
    \right).
\end{equation*}

Следовательно,
%
\begin{equation*}
    \frac{M}{m}
    =
    \frac{2}{1+e^{\chi a}}.
\end{equation*}

Отсюда
%
\begin{equation*}
    1+e^{\chi a}
    =
    \frac{2m}{M},
\end{equation*}
%
и
%
\begin{equation*}
    \rboxed{
        e^{\chi a}
        =
        \frac{2m}{M}-1
    }.
\end{equation*}

Или
%
\begin{equation*}
    \rboxed{
        \chi a
        =
        \ln
        \left(
            \frac{2m}{M}-1
        \right)
    }.
\end{equation*}

Для локализованного состояния необходимо
%
\begin{equation*}
    \chi>0.
\end{equation*}

Значит,
%
\begin{equation*}
    e^{\chi a}>1.
\end{equation*}

Следовательно,
%
\begin{equation*}
    \frac{2m}{M}-1>1,
\end{equation*}
%
откуда
%
\begin{equation*}
    \rboxed{
        M<m
    }.
\end{equation*}

Таким образом, локализованная мода существует,
если примесный атом легче атомов основной цепочки.


% ------------------------------------------------------------
% Частота локализованной моды
% ------------------------------------------------------------

\subsection*{Частота локализованного состояния}

Используем
%
\begin{equation*}
    e^{\chi a}
    =
    \frac{2m-M}{M}.
\end{equation*}

Тогда
%
\begin{align*}
    \omega_{\mathrm{loc}}^{\,2}
    &=
    \frac{\alpha}{m}
    \left(
        e^{\chi a}
        +
        e^{-\chi a}
        +
        2
    \right)
    \\
    &=
    \frac{\alpha}{m}
    \left(
        \frac{2m-M}{M}
        +
        \frac{M}{2m-M}
        +
        2
    \right).
\end{align*}

Приводим выражение в скобках к общему знаменателю:
%
\begin{align*}
    \frac{2m-M}{M}
    +
    \frac{M}{2m-M}
    +
    2
    &=
    \frac{
        (2m-M)^2
        +
        M^2
        +
        2M(2m-M)
    }{
        M(2m-M)
    }
    \\
    &=
    \frac{4m^2}{
        M(2m-M)
    }.
\end{align*}

Следовательно,
%
\begin{equation*}
    \rboxed{
        \omega_{\mathrm{loc}}^{\,2}
        =
        \frac{
            4\alpha m
        }{
            M(2m-M)
        }
    }.
\end{equation*}


% ------------------------------------------------------------
% Положение относительно фононной зоны
% ------------------------------------------------------------

\subsection*{Положение локализованной моды относительно фононной зоны}

Для идеальной цепочки
%
\begin{equation*}
    \omega_{\max}^{\,2}
    =
    \frac{4\alpha}{m}.
\end{equation*}

Рассмотрим отношение
%
\begin{equation*}
    \frac{
        \omega_{\mathrm{loc}}^{\,2}
    }{
        \omega_{\max}^{\,2}
    }
    =
    \frac{
        m^2
    }{
        M(2m-M)
    }.
\end{equation*}

Сравним числитель и знаменатель:
%
\begin{align*}
    m^2
    -
    M(2m-M)
    &=
    m^2
    -
    2mM
    +
    M^2
    \\
    &=
    (m-M)^2.
\end{align*}

При
%
\begin{equation*}
    M\neq m
\end{equation*}
%
имеем
%
\begin{equation*}
    (m-M)^2>0.
\end{equation*}

Следовательно,
%
\begin{equation*}
    \rboxed{
        \omega_{\mathrm{loc}}
        >
        \omega_{\max}
    }.
\end{equation*}

То есть локализованная примесная мода находится
выше полосы частот распространяющихся фононов идеальной цепочки.


% ============================================================
% Задача 3
% ============================================================

\task[3]{Дефект жёсткости в линейной цепочке}

\condition
Все атомы цепочки имеют одинаковую массу $m$,
но жёсткость одной из пружин отличается от остальных.
Пусть её жёсткость равна $\alpha_0$,
тогда как все остальные связи имеют жёсткость $\alpha$.
Определить условие существования локализованного состояния.


% ------------------------------------------------------------
% РИСУНОК 4
% ------------------------------------------------------------

\begin{rbox}[left=18pt,right=18pt,top=14pt,bottom=10pt]{[]}
    \centering

    \begin{tikzpicture}[
        >=Latex,
        every node/.style={text=rtext}
    ]

        % Обычная левая пружина
        \draw[
            draw=rtext,
            line width=0.8pt
        ]
        plot[
            domain=0:1,
            samples=70,
            variable=\t
        ]
        (
            {-3.72+1.44*\t},
            {0.11*sin(1440*\t)}
        );

        % Обычная пружина слева от дефекта
        \draw[
            draw=rtext,
            line width=0.8pt
        ]
        plot[
            domain=0:1,
            samples=70,
            variable=\t
        ]
        (
            {-1.72+1.44*\t},
            {0.11*sin(1440*\t)}
        );

        % Дефектная пружина
        \draw[
            draw=rc3,
            line width=1.3pt
        ]
        plot[
            domain=0:1,
            samples=70,
            variable=\t
        ]
        (
            {0.28+1.44*\t},
            {0.13*sin(1440*\t)}
        );

        % Обычная правая пружина
        \draw[
            draw=rtext,
            line width=0.8pt
        ]
        plot[
            domain=0:1,
            samples=70,
            variable=\t
        ]
        (
            {2.28+1.44*\t},
            {0.11*sin(1440*\t)}
        );

        % Края
        \node at (-5.0,0) {$\ldots$};
        \node at (5.0,0) {$\ldots$};

        % Атомы
        \foreach \x in {-4,-2,0,2,4}{
            \filldraw[
                fill=rbgsecondary,
                draw=rtext,
                line width=0.9pt
            ]
            (\x,0) circle (0.28);

            \node[above=5pt] at (\x,0) {$m$};
        }

        % Номера около дефекта
        \node[below=5pt] at (0,0) {$0$};
        \node[below=5pt] at (2,0) {$1$};

        % Жёсткости
        \node[above] at (-3,0.18) {$\alpha$};
        \node[above] at (-1,0.18) {$\alpha$};
        \node[above,text=rc3] at (1,0.20) {$\alpha_0$};
        \node[above] at (3,0.18) {$\alpha$};

    \end{tikzpicture}

    \captionsetup{
        font=footnotesize,
        labelfont=bf,
        skip=5pt
    }

    \captionof{figure}{
        Одномерная цепочка одинаковых атомов массы $m$
        с одной дефектной связью.
        Обычные пружины имеют жёсткость $\alpha$,
        а пружина между атомами $0$ и $1$ ---
        жёсткость $\alpha_0$.
    }

    \label{fig:spring_impurity}
\end{rbox}


Для атомов, не соседствующих с дефектной пружиной,
уравнение движения остаётся обычным:
%
\begin{equation*}
    \rboxed{
        m\ddot u_n
        =
        \alpha
        \left(
            u_{n+1}
            +
            u_{n-1}
            -
            2u_n
        \right)
    },
    \qquad
    n\neq0,1.
\end{equation*}

Для атома $n=0$:
%
\begin{equation*}
    \rboxed{
        m\ddot u_0
        =
        \alpha(u_{-1}-u_0)
        +
        \alpha_0(u_1-u_0)
    }.
\end{equation*}

Для атома $n=1$:
%
\begin{equation*}
    \rboxed{
        m\ddot u_1
        =
        \alpha_0(u_0-u_1)
        +
        \alpha(u_2-u_1)
    }.
\end{equation*}


% ------------------------------------------------------------
% Затухающая мода
% ------------------------------------------------------------

\subsection*{Затухающая мода}

Как и в предыдущей задаче,
для локализованного состояния вне дефекта берём решение,
амплитуда которого экспоненциально уменьшается.

Обозначим
%
\begin{equation*}
    q=e^{-\chi a},
    \qquad
    0<q<1.
\end{equation*}

Тогда для ближайших к дефекту атомов
%
\begin{equation*}
    u_{-1}
    =
    -q\,u_0,
    \qquad
    u_2
    =
    -q\,u_1.
\end{equation*}

Знак минус снова соответствует чередованию знака
смещений соседних атомов.

Вдали от дефекта получаем то же соотношение,
что и в предыдущей задаче:
%
\begin{equation*}
    \rboxed{
        m\omega^2
        =
        \alpha
        \left(
            e^{\chi a}
            +
            e^{-\chi a}
            +
            2
        \right)
    }.
\end{equation*}


% ------------------------------------------------------------
% Граничные уравнения
% ------------------------------------------------------------

\subsection*{Уравнения для атомов около дефекта}

Положим
%
\begin{equation*}
    u_0(t)
    =
    U_0e^{-i\omega t},
    \qquad
    u_1(t)
    =
    U_1e^{-i\omega t}.
\end{equation*}

Для атома $0$ имеем
%
\begin{equation*}
    -m\omega^2U_0
    =
    \alpha(-qU_0-U_0)
    +
    \alpha_0(U_1-U_0).
\end{equation*}

После переноса знаков:
%
\begin{equation*}
    m\omega^2U_0
    =
    \alpha(1+q)U_0
    +
    \alpha_0(U_0-U_1).
\end{equation*}

Аналогично для атома $1$:
%
\begin{equation*}
    m\omega^2U_1
    =
    \alpha(1+q)U_1
    +
    \alpha_0(U_1-U_0).
\end{equation*}

Получаем систему
%
\begin{equation*}
    \rboxed{
    \begin{aligned}
        m\omega^2U_0
        &=
        \alpha(1+q)U_0
        +
        \alpha_0(U_0-U_1),
        \\
        m\omega^2U_1
        &=
        \alpha(1+q)U_1
        +
        \alpha_0(U_1-U_0).
    \end{aligned}
    }
\end{equation*}


% ------------------------------------------------------------
% Симметричная и антисимметричная моды
% ------------------------------------------------------------

\subsection*{Симметрия относительно дефектной связи}

Система симметрична относительно середины дефектной пружины,
поэтому возможны два типа решения.

Для симметричной моды
%
\begin{equation*}
    U_1=U_0.
\end{equation*}

Тогда дефектная пружина не растягивается:
%
\begin{equation*}
    U_1-U_0=0,
\end{equation*}
%
и параметр $\alpha_0$ вообще выпадает из уравнения.

Такое решение не даёт искомой локализованной моды,
связанной с дефектной пружиной.

Поэтому рассматриваем антисимметричную моду:
%
\begin{equation*}
    \rboxed{
        U_1=-U_0
    }.
\end{equation*}

Тогда
%
\begin{equation*}
    U_0-U_1
    =
    2U_0.
\end{equation*}

Первое граничное уравнение принимает вид
%
\begin{equation*}
    \rboxed{
        m\omega^2
        =
        \alpha
        \left(
            1+e^{-\chi a}
        \right)
        +
        2\alpha_0
    }.
\end{equation*}


% ------------------------------------------------------------
% Условие на chi
% ------------------------------------------------------------

\subsection*{Условие локализации}

С другой стороны, вдали от дефекта
%
\begin{equation*}
    m\omega^2
    =
    \alpha
    \left(
        e^{\chi a}
        +
        e^{-\chi a}
        +
        2
    \right).
\end{equation*}

Приравниваем два выражения:
%
\begin{equation*}
    \alpha
    \left(
        e^{\chi a}
        +
        e^{-\chi a}
        +
        2
    \right)
    =
    \alpha
    \left(
        1+e^{-\chi a}
    \right)
    +
    2\alpha_0.
\end{equation*}

Сокращаем одинаковые слагаемые:
%
\begin{equation*}
    \alpha
    \left(
        e^{\chi a}+1
    \right)
    =
    2\alpha_0.
\end{equation*}

Следовательно,
%
\begin{equation*}
    \rboxed{
        e^{\chi a}
        =
        \frac{2\alpha_0}{\alpha}-1
    }.
\end{equation*}

Или
%
\begin{equation*}
    \rboxed{
        \chi a
        =
        \ln
        \left(
            \frac{2\alpha_0}{\alpha}-1
        \right)
    }.
\end{equation*}

Для локализованного состояния необходимо
%
\begin{equation*}
    \chi>0.
\end{equation*}

Поэтому
%
\begin{equation*}
    e^{\chi a}>1.
\end{equation*}

Следовательно,
%
\begin{equation*}
    \frac{2\alpha_0}{\alpha}-1>1,
\end{equation*}
%
откуда
%
\begin{equation*}
    \rboxed{
        \alpha_0>\alpha
    }.
\end{equation*}

Таким образом, локализованное состояние возникает,
если дефектная связь жёстче остальных связей цепочки.


% ------------------------------------------------------------
% Частота локализованной моды
% ------------------------------------------------------------

\subsection*{Частота локализованной моды}

Введём обозначение
%
\begin{equation*}
    x
    =
    \frac{\alpha_0}{\alpha}.
\end{equation*}

Тогда
%
\begin{equation*}
    e^{\chi a}
    =
    2x-1.
\end{equation*}

Используем выражение для частоты:
%
\begin{equation*}
    \omega_{\mathrm{loc}}^{\,2}
    =
    \frac{\alpha}{m}
    \left(
        e^{\chi a}
        +
        e^{-\chi a}
        +
        2
    \right).
\end{equation*}

Подставляем
%
\begin{equation*}
    e^{\chi a}=2x-1,
    \qquad
    e^{-\chi a}
    =
    \frac{1}{2x-1}.
\end{equation*}

Получаем
%
\begin{equation*}
    \omega_{\mathrm{loc}}^{\,2}
    =
    \frac{\alpha}{m}
    \left(
        2x-1
        +
        \frac{1}{2x-1}
        +
        2
    \right).
\end{equation*}

Приводим к общему знаменателю:
%
\begin{align*}
    2x-1
    +
    \frac{1}{2x-1}
    +
    2
    &=
    2x+1
    +
    \frac{1}{2x-1}
    \\
    &=
    \frac{
        (2x+1)(2x-1)+1
    }{
        2x-1
    }
    \\
    &=
    \frac{4x^2}{2x-1}.
\end{align*}

Следовательно,
%
\begin{equation*}
    \rboxed{
        \omega_{\mathrm{loc}}^{\,2}
        =
        \frac{4\alpha}{m}
        \frac{x^2}{2x-1}
    },
    \qquad
    x=\frac{\alpha_0}{\alpha}.
\end{equation*}


% ------------------------------------------------------------
% Сравнение с верхней границей зоны
% ------------------------------------------------------------

\subsection*{Положение связанного состояния}

Для идеальной цепочки
%
\begin{equation*}
    \omega_{\max}^{\,2}
    =
    \frac{4\alpha}{m}.
\end{equation*}

Поэтому
%
\begin{equation*}
    \frac{
        \omega_{\mathrm{loc}}^{\,2}
    }{
        \omega_{\max}^{\,2}
    }
    =
    \frac{x^2}{2x-1}.
\end{equation*}

При существовании локализованного состояния
%
\begin{equation*}
    x>1.
\end{equation*}

Сравним это отношение с единицей:
%
\begin{equation*}
    \frac{x^2}{2x-1}>1
    \quad\Longleftrightarrow\quad
    x^2>2x-1.
\end{equation*}

Переносим всё в левую часть:
%
\begin{equation*}
    x^2-2x+1>0.
\end{equation*}

То есть
%
\begin{equation*}
    (x-1)^2>0.
\end{equation*}

Для дефектной связи
%
\begin{equation*}
    x\neq1,
\end{equation*}
%
поэтому
%
\begin{equation*}
    \rboxed{
        \omega_{\mathrm{loc}}
        >
        \omega_{\max}
    }.
\end{equation*}

Следовательно, локализованное состояние,
если оно существует, лежит выше разрешённой фононной полосы
идеальной одномерной цепочки.

\end{document}